% Generated draft: requires mathematical and source review.
\documentclass[11pt]{article}
\usepackage[margin=25mm]{geometry}
\usepackage{fontspec}
\setmainfont{Latin Modern Roman}
\usepackage{amsmath,amssymb,mathtools,unicode-math}
\setmathfont{Latin Modern Math}
\usepackage{enumitem,needspace,xurl,hyperref}
\setlength{\parindent}{0pt}
\setlength{\parskip}{5pt}
\setlength{\emergencystretch}{3em}
\newcommand{\N}{\mathbb{N}}
\newcommand{\Z}{\mathbb{Z}}
\newcommand{\Q}{\mathbb{Q}}
\newcommand{\R}{\mathbb{R}}
\newcommand{\C}{\mathbb{C}}
\newcommand{\F}{\mathbb{F}}
\newcommand{\A}{\mathbb{A}}
\newcommand{\PP}{\mathbb{P}}
\newcommand{\E}{\mathbb{E}}
\newcommand{\eps}{\varepsilon}
\newcommand{\dd}{\,\mathrm{d}}
\begin{document}
\section*{UnsolvedMath ID: 11100025}
\textbf{Source problem code:} catalogue entry
\subsection{Definitions and mathematical statement}
A $C^*$-algebra is a complex Banach $*$-algebra satisfying $\|a^*a\|=\|a\|^2$. Kasparov's $KK(A,B)$ is a bivariant group, contravariant in $A$ and covariant in $B$, in its intended class of $C^*$-algebras. The excerpt does not specify separability, unitality, or other analytic hypotheses.

\textbf{Applications problem 1 (Hovey).} Develop a useful homotopy theory for a suitable category of $C^*$-algebraic objects by constructing a model structure and then a stabilization. Seek a stable homotopy category $\mathcal D$ and an appropriate comparison $j$ for which
\[
 KK(A,B)\cong\operatorname{Hom}_{\mathcal D}(jA,jB)
\]
naturally for the intended algebras. Also seek representability of $C^*$-algebra $K$-theory. In a $KK$-realizing interpretation, its degree-zero instance is
\[
 K_0(A)\cong\operatorname{Hom}_{\mathcal D}(j\C,jA),
\]
with higher grading expressed using the specified suspension convention. The source suggests these goals may be the same, but does not establish their equivalence or specify the category, weak equivalences, stabilization, or representing objects.

Hovey proposes allowing pro-$C^*$-algebraic objects to accommodate inverse-limit constructions, citing Phillips. A pro-$C^*$-algebra in that framework is an appropriate topological projective limit of $C^*$-algebras. The source's blanket assertion that $C^*$-algebras are not closed under inverse limits requires a category/limit convention: bounded $C^*$-products and unrestricted topological inverse-limit constructions differ. This motivation must not be treated as an unqualified categorical nonexistence theorem.

He further suggests that a model category may be fibrantly rather than cofibrantly generated because mapping into these objects seems easier than mapping out. This is a proposed direction, not an imposed theorem or an established model structure. The broader aim, compared with Voevodsky's algebraic-geometric homotopy theory, is to allow theories without an axiom forcing Bott periodicity. There is a formulation issue to resolve: a category realizing $KK$ as all the displayed morphisms has its built-in periodicity, while a broader nonperiodic homotopy theory would need a specified relationship to that periodic part, for example a subsequent comparison or localization. The source does not supply this relationship and no such construction is asserted here.
\subsection{Short English statement}
Build a model-category and stable-homotopy framework for suitable $C^*$-algebraic objects, realize $KK$ as morphisms and $K$-theory representably, and clarify how this periodic theory fits into the proposed broader nonperiodic setting.
\subsection{Sources}
\begin{itemize}
\item Local UnsolvedMath registry, permanent ID 11100025, source code catalogue entry; attributed to Mark Hovey, no proposal year supplied.
\item \emph{Hovey's Algebraic Topology Problem List}, Applications, item 1; surviving mirror \url{https://www-users.cse.umn.edu/~tlawson/hovey/}. The registry records HTML-ordered-list extraction and access on 2026-07-29; the mirror contains 65 top-level items on nine category pages.
\item N. C. Phillips, ``Inverse limits of $C^*$-algebras,'' \emph{Journal of Operator Theory} \textbf{19} (1988), no. 1, 159--195; MR 90c:46090, exactly as supplied in the source.
\item The source names Claude Schochet as knowledgeable in the area, says Jim McClure shares the viewpoint, and tentatively mentions Marius Dadarlat; these are Hovey's attributed remarks, not claims of a specific construction by those individuals. Voevodsky is named as an analogy. The summary's modern categorical $KK$ assertion has no precise reference here.
\end{itemize}
\end{document}
